% mcm-paper source binding (non-rendering metadata)
% Repository HEAD: 66f49838fe3060c9c0286cb5aa87613ac20b9fc3
% Validation report source: 代码/模型检验与最终结果报告.md @ a3970ef9eb51b87bd1c564a98d2cdfef213373a4
% Q3 figure source: a3970ef9eb51b87bd1c564a98d2cdfef213373a4

\section{模型检验}\label{sec:validation}

为检验结果可靠性，本文从独立输出验证、关键目标认证及候选池参数敏感性三个方面进行复核。三问正式方案均由独立程序重新读取交付CSV并重算约束与指标，避免仅依赖求解器内部状态。

\subsection{模型有效性与结果验证}\label{subsec:validation-results}

\subsubsection{独立验证与认证边界}

独立验证覆盖人员完整性、路线闭环、动态容量、燃油与加油等公共约束；问题三进一步检查人员时间窗、具体飞机兼容、任务冲突及30 min周转。三问正式结果及认证状态汇总于表~\ref{tab:validation-summary}。

\begin{table}[htbp]
\centering
\caption{三问正式结果验证与认证状态}
\label{tab:validation-summary}
\small
\begin{tabular}{c c c c}
\toprule
对象 & 服务结果 & 独立验证 & 认证说明\\
\midrule
问题一 & 1600人全部完成 & 通过 & 当前最好已验证可行解\\
问题二 & 4000人全部完成 & 通过 & 当前最好已验证可行解\\
问题三 & 3840+160人全部完成 & 通过 & 完整问题三未全局认证\\
问题三阶段B & 临时160/160 & 通过 & 受限模型内最优，相对间隙为0\\
\bottomrule
\end{tabular}
\end{table}

其中问题三阶段B在冻结的$B_0=35114$ min和当前有限时间可行架次列池下，临时人员目标的当前最好可行值与求解器最优界均为160，因此

\[
N_{\mathrm{temp},\Omega}^{*}=160.
\tag{9-1}\label{eq:q3-stageb-certification}
\]

160同时等于临时人员总数，故该目标达到人数意义上的理论绝对上限。该认证只针对阶段B的接纳目标；问题一、问题二以及问题三阶段A、阶段C仍按经验证可行结果解释，不扩大为完整路线空间的全局最优。

\subsection{稳定性与敏感性分析}\label{subsec:sensitivity}

\subsubsection{问题一候选加油参数敏感性}

该敏感性实验采用独立且固定的求解设置，仅用于比较不同加油候选保留数对结果的影响，并非对14759 min正式方案的重复求解。在其余设置保持一致的条件下，将加油候选保留数由4调整为8，结果见表~\ref{tab:q1-refuel-sensitivity}。

\begin{table}[htbp]
\centering
\caption{问题一候选加油参数敏感性}
\label{tab:q1-refuel-sensitivity}
\small
\setlength{\tabcolsep}{4pt}
\begin{tabular}{c c c c c c}
\toprule
加油候选保留数 & 路线池 & $T^{air}$/min & $T^{pax}$/人$\cdot$min & 燃油/kg & 架次数\\
\midrule
4 & 1603 & 14837 & 127784 & 119359.5 & 88\\
8 & 1622 & 14908 & 124166 & 119365.4 & 89\\
\bottomrule
\end{tabular}
\end{table}

由表~\ref{tab:q1-refuel-sensitivity}可见，参数由4增至8后，路线池仅扩大1.185\%，飞机时间变化0.479\%，燃油变化不足0.005\%。人员在途时间下降2.831\%，总体表明主要运行指标对该参数较稳定。不同指标的小幅反向变化来自有限时间词典序搜索，不能据此认为更大的候选池必然更差。

\subsubsection{问题三时间列密度敏感性}

问题三在基准采样方案“最早--中点--最晚”的基础上构造加密采样方案，并增加25\%和75\%分位点。图~\ref{fig:q3-time-density}显示，时间可行架次列与服务选项规模分别增加约85\%和94\%；与此同时，加密采样方案的阶段A及阶段C飞机使用时间均下降，燃油由301164.1 kg降至289768.1 kg，座位利用率由44.07\%升至45.52\%，说明更细时间嵌入能够提供更多有效排班组合。

\begin{figure}[htbp]
    \centering
    \includegraphics[width=0.96\textwidth]{figures/14时间列密度敏感性.pdf}
    \caption{问题三时间列密度敏感性分析}
    \label{fig:q3-time-density}
\end{figure}

但在相同限时求解条件下，阶段B当前可行接纳数由143降至112。由于加密采样候选集是在基准方案上扩展得到，可行域并未缩小，因此该差异应归因于模型规模增大后的组合搜索困难，而不能解释为“时间列越密，可服务人数越少”。这一判断还得到正式求解的验证：在更充分搜索后，阶段B最终达到160/160，并获得相对间隙为0的严格认证。因此143和112仅反映有限时间内当前最好可行解的质量，而不是临时人员服务能力上限。

综上，三问正式交付文件均通过独立验证；问题三阶段B的临时人员接纳目标得到严格认证。现有敏感性实验表明，问题一对加油候选保留数总体稳定，而问题三时间列加密能够提升排班灵活性，但会显著增加有限时间搜索负担。独立可行性复核、受限目标认证与参数扰动实验共同构成了本文的模型检验链。
